\section{Why Weak ETH rules out spectral expansion}\label{sec:graphs}

Transposition sends each bit from its input row to its output column.
A short tape computation would implement all these transfers with a
bounded-degree network. To rule this out we will place the vertices
in a metric space where the network edges are short and the
row-to-column distances are typically large. The issue here is the
existence of such placements. Uniform expansion in the relevant
metrics would instead allow a deterministic algorithm to decide
3-SAT with a power saving in the number of variables.

\subsection{Wire costs and expansion}

In a rational tensor network each edge of width $q$ carries a string of
$q$ bits. A vertex table assigns a rational coefficient to each tuple
of incident strings and external indices. To contract the network,
fix the external indices, multiply the table entries, and sum over all
internal strings. A network with no external indices is \emph{closed}.
A \emph{bit network} has width one on every edge.

The metric records the cost of sharing information. In the tensor
language, a wire is equality of two external labels. A one-bit wire
allows the joint labels $(0,0)$ and $(1,1)$, each with coefficient one.
To obtain a metric, take tensors over $\mathbb Q$ at a finite set of at
least two sites. A spectral point in Strassen's asymptotic spectrum
takes nonnegative real values, preserves zero, one, direct sums and
tensor products, and is monotone under a linear map at each site
\cite{Strassen}. Its base-two logarithm, with a positive normalization,
is a \emph{price}. The prices of one-bit wires give a pseudometric on the
sites. Thus a $q$-bit wire between two sites at distance $d$ costs $qd$.

We use the class $\mathcal X$ of complete $\operatorname{CAT}(0)$
spaces carrying compatible spectral prices for rational tensors on
all finite subsets, with their distances given by the wire prices.
Every finite spectral wire pseudometric occurs in a member of
$\mathcal X$ after sites at distance zero are identified. Rescaling
prices and distances together is allowed
\cite[Sections~\ext{G}{sec:spectra} and~\ext{G}{sec:limits},
Theorem~\ext{G}{thm:geometry}]{oaG}.
Place each vertex $u$ at $x_u$ in a member of $\mathcal X$. Applying
the vertex tables as local maps to the product of the edge wires gives
the contraction. Tables at coincident sites can be applied jointly.
It follows by monotonicity and additivity of the price that any product
of equality wires obtained in this way costs at most
\[
                 \sum_{e=uv} q_e\,d(x_u,x_v).
\]
For $r$ copies the right-hand side is multiplied by $r$, even when
the local maps act jointly on all copies. If all vertices coincide,
every wire costs zero.

Characteristic zero \emph{admits ports} when there are symmetric
regular multigraphs of arbitrarily large order, with a common bound
on their degrees, and a constant $C$ such that
\[
                 \mathbb E\,d(x_U,x_V)^2
                      \le C\,\mathbb E\,d(x_I,x_J)^2
\]
for every placement of every such graph in $\mathcal X$. Here $U,V$
are independent uniform vertices and $(I,J)$ is a uniform directed
arc, with a diagonal arc counted once in the degree. This is the
nonlinear Poincar\'e inequality for the wire metrics, which depend
only on the characteristic \cite[Definition~\ext{G}{def:ports}]{oaG}.

\begin{theorem}[Expansion and 3-SAT]\label{thm:noports}
If characteristic zero admits ports, then for some $\delta>0$ a
deterministic algorithm decides 3-SAT on $v$ variables in
$2^{O(v^{1-\delta})}$ bit operations times a polynomial in the input
length. In particular, under variable-count Weak ETH, characteristic
zero does not admit ports.
\end{theorem}

\begin{proof}
Suppose characteristic zero admits ports. By the enumerable contraction
theorem \cite[Theorem~\ext{P}{thm:closed}]{oaP}, there is $0<\eta<1$
such that a closed bit network of volume $s$ (vertices plus edges)
and maximum degree at most $s^{1/20}$ can be contracted
deterministically in $2^{O(s^{1-\eta})}$ bit operations plus the time
for that many local entry queries when its entries are integers of
polynomial bit length. The topology determines an enumerable formula
with fixed rational constants. For the clause and equality tables
below, each query takes constant work given its arguments.

A general 3-CNF can have order $v^3$ clauses. To get a power saving
in $v$, we branch into formulas of smaller volume by sunflower
sparsification \cite{IPZ}. Choose a fixed rational $0<\alpha<1/4$
so small that $(1+4\alpha)(1-\eta)<1$. Put $r=\lceil v^\alpha\rceil$
and $R=r^2$, taking $v\ge2$. Remove tautologies and duplicate literals,
and handle an empty clause directly. Regard each remaining clause as a
set of literals. At the outset and after every branch discard duplicate
clauses and strict supersets of other clauses.

At each node choose the first applicable rule below, breaking ties
canonically. Each rule selects clauses with a common center and
pairwise disjoint, nonempty remainders, the petals.
\begin{enumerate}[leftmargin=*]
\item If a literal occurs in at least $r$ two-literal clauses, take
$r$ of them with that literal as center.
\item If a pair of literals occurs in at least $R$ three-literal
clauses, take $R$ of them with that pair as center.
\item If a literal occurs in more than $2R^2$ three-literal clauses,
choose $R$ of them with that literal as center and disjoint two-literal
petals. A petal excludes at most $2R$ candidates since the preceding
rule does not apply, so greedy selection suffices.
\end{enumerate}
One child adds the center as a clause. The other adds every petal as
a clause. An assignment satisfying the parent either satisfies the
center or must satisfy every petal. Both children imply the parent.
Each new clause is a strict subset of a selected clause, so no
existing clause is contained in it. The new petals are disjoint and
hence all survive simplification.

Consider a root-to-leaf path. At most $2v$ unit clauses are added,
since each persists. The initial run of the first rule adds only units.
After that run, each literal belongs to fewer than $r$ pairs. Subsequent
branches on triples can add one pair or a set of disjoint pairs, so
they raise each pair degree by at most one. The first rule has priority
and only deletes pairs. Thus pair degrees stay at most $r$. At any
node after the initial run there are $O(vr)$ surviving pairs, and
each subsequent unit can delete at most $r$ of them. Only units
delete pairs, so the total number of new pairs along the path is
$O(vr)$. Each step adds a unit or a pair, giving depth $O(vr)$.
Taking the child that adds all petals consumes at least $r$ new
units or $R$ new pairs. Such a choice occurs only $O(v/r)$ times.
Recording the depths of these choices determines the leaf, hence
there are at most $2^{O((v/r)\log v)}$ leaves. Finding the branches
and simplifying at each node take polynomial time in the input length.

At a leaf the three rules all fail, so there are $O(vR^2)$ clauses.
Construct a closed bit network with a clause indicator at each clause
and a tree of ternary equality tables to distribute each variable to
its occurrences. For each unused variable, a wire joining two all-ones unary tables
supplies its two possible values.
Every assignment has a unique compatible labeling of the copy wires,
and the product of the indicators tests all clauses. Its contraction
is therefore the number of satisfying assignments of this leaf.
The network has degree at most three and volume $O(v^{1+4\alpha})$.
For volumes too small to meet the degree hypothesis, direct contraction
costs a constant.

Accept if any leaf contraction is positive. The logarithm of the total
running time, apart from the polynomial input factor, is at most
\[
       O\!\left(v^{1-\alpha}\log v
             + v^{(1+4\alpha)(1-\eta)}\right).
\]
Taking $0<\delta<\min\{\alpha,1-(1+4\alpha)(1-\eta)\}$ gives the
claimed bound, which variable-count Weak ETH excludes.
\end{proof}

\noindent
For the tape lower bound we use the quantitative consequence of
nonexpansion from \cite[Theorem~\ext{P}{thm:loggraph}]{oaP}. Under
Weak ETH, for a
regular graph on $P\ge2$ vertices of fixed positive degree, the
Poincar\'e constant valid for all placements in $\mathcal X$ must
grow at least as $\log P$, up to a degree-dependent factor. One can
therefore choose a placement with mean squared distance one between
independent vertices and mean squared distance $O(1/\log P)$ along
an edge. Lemma~\ref{lem:geometry} applies this fact to the graph of a
tape computation and controls distances to the row and column terminals.
